% HAND-WRITTEN fragment -- NOT generated by maestro2tex. Input from
% AnalogChapter.tex immediately after Pixel.tex, which is generated.
%
% Salvaged from the retired note_calibration_seq.tex step 5 (current zero) and
% the triage table of ~/vestarv/docs/afe_rev2_debug_mux.html §9/§9.1:
% per-channel slot 5 = ce_amp, slot 6 = re_buf (buffered at source), slot 7 =
% we_sum behind a series guard T-gate, slot 8 = we_pad. The discriminating pair
% is CE against WE; both faults read as a stuck code at the converter.
% Numbers: 3 sigma of A1 offset = 15.6 mV (tab_ampab_mc_offset); 0.824 nA/LSB at
% the top gain code (ss:tiarprog); zero-current sigma from sss:pixel-mc.
% Library cell names are kept out of the rendered prose.

\subsubsection{Measuring this block on chip}

Two nodes separate the four faults that all read as a stuck converter code.
Select per-channel slot 5 and confirm the counter-electrode drive sits mid-rail
and follows the reference code; then slot 6, the reference-electrode buffer
output, which must equal the commanded reference within \SI{16}{\milli\volt}
(\(3\sigma\) of A1's offset). A railed counter electrode with the summing node
at the common mode is a control-loop or electrode fault; an unrailed counter
electrode with the summing node off the common mode is a transimpedance fault;
the converter reports a rail for both, which is why this pair and not the
converter is the triage. Read the summing node itself at slot 7, through its
series guard and with the mid node clamped --- it is a virtual ground at
\SI{0.824}{\nano\ampere} per LSB at the top gain code, so an ungranted site must
stay doubly isolated from it. For the zero-current constant, open the
working-electrode isolation switch so the cell current is zero by construction,
read the transimpedance output at every gain code the channel will use, and
subtract the converter's zero code; the uncalibrated spread is \(\sigma =
12.1\)--\SI{167.3}{LSB} over the four taps (Section~\ref{sss:pixel-mc}). This
test fabric is design intent; it is not yet in the schematics.
